\chapter{Financial Generative Modelling}
\label{ch:financial}

\begin{tcolorbox}[feynbox={Sky}{BBg},title={Before and After}]
\textbf{Before this chapter you need:} log returns (\S\ref{sec:logreturns}),
expectation and variance (\S\ref{sec:expectation}--\S\ref{sec:variance}),
covariance (\S\ref{sec:covariance}), quantiles and CDFs
(\S\ref{sec:distributions}), stationarity and GARCH
(\S\ref{sec:stationarity}, \S\ref{sec:garch}), tail index
(\S\ref{sec:moments}).\\
\textbf{After it you will understand:} the seven stylised facts with this
project's measured values, what VaR and expected shortfall report, why the Sharpe
ratio is misleading on heavy-tailed data, and what a structural break does to
every assumption in the preceding chapters --- with the MOEX 2022 suspension as
the worked case.
\end{tcolorbox}

This chapter is the domain half of the problem. It states what a synthetic return
series must reproduce to be useful, and it does so in measurements rather than
adjectives.

%% ─────────────────────────────────────────────────────────────────────────────
\section{Returns, Volatility and Scaling}
\label{sec:volatility}

\situation{
Two funds report their year. One says ``up \$4 million''; the other ``up 12\%''.
Only the second can be compared to anything. And neither says how nervous the year
was --- a steady 12\% and a 12\% arrived at via two crashes are different products.
}

\problem{
Level, change and risk are three different quantities, and financial data supplies
only the first directly. Everything else is derived, and the derivations carry
conventions worth stating explicitly.
}

\workedarith{
\textbf{Return, restated.} From \S\ref{sec:logreturns},
$r_t = \ln(P_t/P_{t-1})$, additive across time.

\medskip
\textbf{Volatility is a standard deviation of returns.} For the five markets in
this project, measured over 2006--2026:

\medskip
\begin{tabular}{@{}lrr@{}}
\toprule
Market & Daily $\sigma$ & Annualised \\
\midrule
MOEX     & $1.87\%$ & $29.7\%$ \\
BOVESPA  & $1.63\%$ & $25.9\%$ \\
\bottomrule
\end{tabular}

\medskip
\textbf{The $\sqrt{252}$ rule, derived.} If daily returns were independent with
common variance $\sigma^2$, then over $T$ days variances add
(\S\ref{sec:variance}):
\[
  \Var\!\left(\sum_{t=1}^{T} r_t\right) = T\sigma^2 ,
  \qquad\text{so}\qquad
  \sigma_T = \sigma\sqrt{T} .
\]
With 252 trading days, $\sqrt{252} = 15.87$. For MOEX:
$1.87\% \times 15.87 = 29.7\%$. \checkmark

\medskip
\textbf{Why the rule is wrong here.} It assumed independence. Returns are
uncorrelated but not independent (\S\ref{sec:covariance}), and volatility clusters,
so realised annual volatility varies far more across years than the formula
implies. The convention is retained for comparability, not accuracy.

\medskip
\textbf{Realised volatility.} An alternative estimate over a window:
\[
  \mathrm{RV}_t = \sqrt{\frac{1}{n}\sum_{i=t-n+1}^{t} r_i^2} ,
\]
which is what the rolling-volatility panel of the stylised-facts figure plots.
}

\formaldef{
For a price series $\{P_t\}$:
\[
  r_t = \ln\frac{P_t}{P_{t-1}} , \qquad
  \sigma = \sqrt{\Var(r_t)} , \qquad
  \sigma_{\text{annual}} = \sigma\sqrt{252} .
\]
\textbf{Unconditional volatility} $\sigma$ is the long-run standard deviation;
\textbf{conditional volatility} $\sigma_t$ is the standard deviation given
information up to $t-1$, which is what GARCH models (\S\ref{sec:garch}).

The distinction is central: $\sigma$ is a single number for the whole history,
$\sigma_t$ is a series.
}

\analogybreak{
Volatility is a symmetric measure applied to an asymmetric phenomenon. It treats a
$3\%$ rise and a $3\%$ fall identically, whereas an investor does not, and it says
nothing about the shape of the tail --- two series can share a standard deviation
while one has a tail index of $2.5$ and the other is Gaussian. That is exactly why
\texttt{std\_diff} alone is a weak metric and \texttt{tail\_index\_diff} is
computed alongside it.
}

\whyhere{
Volatility is the scale against which every guard threshold in the pipeline is
set. The post-generation check requires a generated standard deviation between
$0.001$ and $0.15$, bracketing the $0.9$--$2.0\%$ daily range measured across the
five markets --- loose enough to admit any plausible market, tight enough to catch
the collapse of \S\ref{sec:modecollapse} at $0.0002$.
}

%% ─────────────────────────────────────────────────────────────────────────────
\section{The Seven Stylised Facts}
\label{sec:stylisedfacts}

\situation{
A forger studying banknotes lists the features to reproduce: the watermark, the
thread, the paper weight, the microprint. The list is what makes forgery
checkable --- and equally, what makes detection possible.
}

\problem{
``Realistic returns'' must become a finite list of measurable properties, or there
is no way to score a generator. Cont (2001) assembled the canonical list; this
section attaches this project's own measurements to it.
}

\workedarith{
\medskip
\begin{tabular}{@{}rlll@{}}
\toprule
\# & Fact & Measured here & Metric \\
\midrule
1 & Heavy tails & $\hat\alpha = 2.54$--$3.21$ & \texttt{tail\_index\_diff} \\
2 & No linear autocorrelation & $\hat\rho(r,1) = +0.03$ & \texttt{acf\_returns\_mae} \\
3 & Volatility clustering & $\hat\rho(|r|,1) = +0.19$ & \texttt{acf\_absolute\_mae} \\
4 & Long memory in volatility & $H = 0.57$--$0.65$ & \texttt{hurst\_diff} \\
5 & Aggregational Gaussianity & --- & not scored \\
6 & ARCH effects & $\hat\rho(r^2,1) = +0.17$ & \texttt{acf\_squared\_mae} \\
7 & Conditional heavy tails & residual kurtosis $> 0$ & \texttt{resid\_kurtosis\_diff} \\
\bottomrule
\end{tabular}

\medskip
\textbf{Fact 1, in context.} A tail index near 3 means
$\Prob(|r| > x) \sim x^{-3}$. The consequence, from \S\ref{sec:moments}: the
fourth moment is infinite in all five markets, so kurtosis is not estimable.

\medskip
\textbf{Facts 2 and 3 together.} This is the pairing that defines the problem.
Direction is unpredictable ($+0.03$, at the edge of the $\pm0.028$ band) while
magnitude is strongly predictable ($+0.19$, seven times the band). A generator
must reproduce \emph{both} --- and producing independent noise satisfies fact 2
perfectly while failing fact 3 completely.

\medskip
\textbf{Fact 5, and why it is not scored.} Aggregational Gaussianity says that as
returns are aggregated over longer horizons, their distribution approaches a
normal --- monthly returns are closer to Gaussian than daily ones. It is a
consequence of the CLT (\S\ref{sec:clt}) applying to sums. This project evaluates
at the daily horizon only, so the fact is real but not measured here; scoring it
would require re-aggregating the synthetic series, which is a natural extension.

\medskip
\textbf{Fact 7, the subtle one.} After dividing out GARCH volatility, the
standardised residuals \emph{still} have excess kurtosis (Bollerslev 1987). Fat
tails are not merely a by-product of clustering --- there is genuine conditional
heavy-tailedness on top. A generator can reproduce facts 1 and 3 and still fail
this, and \S\ref{sec:qmle} explains the QMLE design that makes the test
non-circular.

\medskip
\textbf{The clustering ratio, restated.} From \S\ref{sec:conditional}, MOEX:
$\Prob(A) = 3.8\%$ against $\Prob(A \mid B) = 22.9\%$, a ratio of $6.0$. Facts 3,
4 and 6 are three different instruments measuring that one phenomenon.
}

\formaldef{
The \textbf{stylised facts} of asset returns (Cont 2001) are empirical regularities
observed across markets, instruments and periods. They are descriptive rather than
derived from theory, and their value is that they are stable enough to serve as
falsifiable targets.

A generative model for financial returns is judged by how many it reproduces and
how closely --- which is precisely the evaluation framework of
Chapter~\ref{ch:evaluation}.
}

\analogybreak{
The list is a necessary condition, never a sufficient one. A generator could
reproduce all seven and still be useless --- for instance by memorising and
replaying the training data, which matches every fact perfectly while generating
nothing new. The shuffled-real control makes exactly this point in reverse: it
matches every permutation-invariant fact by construction while destroying every
temporal one.
}

\whyhere{
These seven facts are what the fourteen ranked metrics measure. The mapping is not
one-to-one --- facts 3, 4 and 6 are covered by four metrics between them, and fact
5 by none --- and stating the mapping explicitly is what allows the evaluation to
claim coverage rather than assume it.
}

%% ─────────────────────────────────────────────────────────────────────────────
\section{Value at Risk}
\label{sec:var}

\situation{
A bank asks: on a bad day, how much might we lose? Not the worst conceivable
loss --- that is unbounded and useless for planning --- but a threshold exceeded
only rarely.
}

\problem{
Risk must become a single reportable number for capital rules and internal limits.
Any such number discards information, and knowing exactly which information is
discarded is the difference between using it and being misled by it.
}

\workedarith{
\textbf{Definition by counting.} The $5\%$ VaR is the loss exceeded on $5\%$ of
days --- the $5$th percentile of the return distribution.

\medskip
\textbf{Under a normal distribution.} With $\mu \approx 0$ and
$\sigma = 1.63\%$ (BOVESPA), the $5\%$ quantile of a standard normal is
$z_{0.05} = -1.645$:
\[
  \mathrm{VaR}_{0.05} = -1.645 \times 1.63\% = -2.68\% .
\]
On a \$1M position, about \$26{,}800.

\medskip
\textbf{Why that understates the risk.} The Gaussian assumption is wrong in
exactly the region VaR cares about. At $\hat\alpha \approx 3$ the empirical $5\%$
quantile is further out than the normal predicts, and the gap widens at more
extreme levels: at $1\%$, the Gaussian says $-2.326\sigma$ while a heavy-tailed
distribution says considerably more.

\medskip
\textbf{Empirical VaR instead.} Sort the returns and take the $5$th percentile
directly. With $n = 496$ test observations, that is the $25$th smallest value
($0.05 \times 496 = 24.8$). No distributional assumption needed --- and no ability
to extrapolate beyond the sample either.

\medskip
\textbf{Counting violations.} If the model is correct, exceedances should occur on
$5\%$ of days. Over $496$ days, expect $24.8$. Observing $40$ is $61\%$ too many;
whether that is evidence of a broken model is the question
Chapter~\ref{ch:downstream} answers.
}

\formaldef{
\textbf{Value at Risk} at level $\alpha$ is the $\alpha$-quantile of the return
distribution:
\[
  \mathrm{VaR}_\alpha = F^{-1}(\alpha) ,
  \qquad\text{so}\qquad
  \Prob\bigl(r_t < \mathrm{VaR}_\alpha\bigr) = \alpha .
\]
\textbf{Unconditional} VaR uses the whole-sample distribution and is a single
number. \textbf{Conditional} VaR uses the distribution given information to
$t-1$, typically $\mathrm{VaR}_{\alpha,t} = \hat\mu + \hat\sigma_t z_\alpha$ with
$\hat\sigma_t$ from GARCH, and is a series that moves with the market.
}

\analogybreak{
VaR reports a threshold and says nothing whatever about what lies beyond it. Two
portfolios can share a $5\%$ VaR of $-2.7\%$ while one loses at most $3\%$ on its
worst day and the other loses $40\%$. VaR is also not subadditive --- the VaR of a
combined portfolio can exceed the sum of the parts --- so it can penalise
diversification, which is why expected shortfall (\S\ref{sec:es}) has largely
displaced it in regulation.
}

\whyhere{
VaR is the downstream task in this thesis: a generator is useful if a risk model
fitted to its synthetic output produces well-calibrated VaR forecasts on real
data. That is the TSTR protocol of Chapter~\ref{ch:downstream}. The reference
figure is the real-GARCH coverage error of $0.0259$.
}

%% ─────────────────────────────────────────────────────────────────────────────
\section{Expected Shortfall}
\label{sec:es}

\situation{
Told a river floods above 3 metres one year in twenty, you still do not know
whether a flood year means 3.1 metres or 9. The threshold is not the damage.
}

\problem{
VaR's blindness beyond the threshold is its central defect, and on heavy-tailed
data that is where the loss lives. A measure that averages over the tail rather
than pointing at its edge answers the question actually being asked.
}

\workedarith{
Ten worst days from a sample, as percentage returns:
\[
  -2.8,\ -3.1,\ -3.4,\ -3.9,\ -4.2,\ -5.1,\ -6.0,\ -7.8,\ -11.2,\ -18.6 .
\]

Suppose these are the worst $10$ of $200$ days, so they are the worst $5\%$.

\medskip
\textbf{VaR at $5\%$} is the threshold: $-2.8\%$, the least bad of the ten.

\medskip
\textbf{Expected shortfall} is their average:
\[
  \mathrm{ES} = \frac{-2.8 -3.1 -3.4 -3.9 -4.2 -5.1 -6.0 -7.8 -11.2 -18.6}{10}
  = \frac{-66.1}{10} = -6.61\% .
\]

\medskip
\textbf{A factor of 2.4.} ES is more than twice VaR, and the gap is driven by the
$-18.6\%$ day. Under a Gaussian the ratio would be about $1.25$; the heavier the
tail, the wider the gap.

\medskip
\textbf{Why ES is subadditive and VaR is not.} ES averages over a tail region, and
averages of sums are sums of averages, so combining portfolios cannot increase it
beyond the sum. VaR reads off a single quantile, and quantiles do not add ---
which permits the pathological cases where diversifying raises reported risk.
}

\formaldef{
\textbf{Expected shortfall} (also conditional VaR, CVaR) at level $\alpha$ is
\[
  \mathrm{ES}_\alpha = \E\bigl[r \mid r \leq \mathrm{VaR}_\alpha\bigr]
  = \frac{1}{\alpha}\int_0^{\alpha} F^{-1}(u)\,du .
\]
It is a \textbf{coherent} risk measure: monotone, translation-equivariant,
positively homogeneous and subadditive. VaR fails subadditivity.

Basel III replaced VaR with $97.5\%$ ES for market-risk capital, largely on these
grounds.
}

\analogybreak{
ES is harder to estimate than VaR precisely because it depends on the extreme
tail, where data is scarcest. It also requires a finite mean of the tail, which by
\S\ref{sec:moments} needs $\alpha > 1$ --- comfortably satisfied here at
$2.54$--$3.21$, but not a free assumption. And backtesting ES is genuinely harder:
it is not \emph{elicitable} on its own, so there is no scoring function whose
minimiser is ES, which is why the Kupiec and Christoffersen tests of
Chapter~\ref{ch:downstream} target VaR.
}

\whyhere{
This thesis backtests VaR rather than ES for exactly that reason: Kupiec and
Christoffersen are well-established VaR tests with clear nulls, and the equivalent
machinery for ES is less settled. ES is the more defensible risk measure and the
less tractable evaluation target, and the choice made here favours testability.
}

%% ─────────────────────────────────────────────────────────────────────────────
\section{The Sharpe Ratio and Portfolio Mathematics}
\label{sec:sharpe}

\situation{
Two funds returned 15\% last year. One did it steadily; the other was up 60\% by
June and down again by December. Comparing them on return alone is meaningless ---
the second took far more risk for the same result.
}

\problem{
Performance must be reported per unit of risk, not in isolation. The standard
device does this by dividing by volatility --- which imports every limitation of
volatility as a risk measure, and on heavy-tailed data that is a serious import.
}

\workedarith{
\textbf{The ratio.} With mean excess return $\bar{r} - r_f$ and volatility
$\sigma$:
\[
  S = \frac{\bar{r} - r_f}{\sigma} .
\]

Annual return $12\%$, risk-free $2\%$, volatility $18\%$:
\[
  S = \frac{0.12 - 0.02}{0.18} = 0.556 .
\]

\medskip
\textbf{Annualising from daily.} A daily Sharpe of $0.035$ scales by $\sqrt{252}$
--- the mean scales with $T$ and the volatility with $\sqrt{T}$, so the ratio
scales with $\sqrt{T}$:
\[
  0.035 \times 15.87 = 0.556 .
\]

\medskip
\textbf{Portfolio risk depends on covariance.} Two assets, weights $w_1, w_2$:
\[
  \sigma_p^2 = w_1^2\sigma_1^2 + w_2^2\sigma_2^2 + 2w_1w_2\rho\sigma_1\sigma_2 .
\]
With $\sigma_1 = \sigma_2 = 1.5\%$ and equal weights:

\medskip
\begin{tabular}{@{}lrl@{}}
\toprule
$\rho$ & $\sigma_p$ & \\
\midrule
$1.0$  & $1.50\%$ & no diversification \\
$0.5$  & $1.30\%$ & partial \\
$0.0$  & $1.06\%$ & full, $= 1.5/\sqrt{2}$ \\
$-1.0$ & $0.00\%$ & perfect hedge \\
\bottomrule
\end{tabular}

\medskip
The same two marginals give portfolio risk anywhere in a $42\%$ range depending
only on $\rho$ --- the point already made in \S\ref{sec:multivariate}.
}

\formaldef{
The \textbf{Sharpe ratio} is
\[
  S = \frac{\E[r] - r_f}{\sqrt{\Var(r)}} ,
\]
scaling as $\sqrt{T}$ under aggregation when returns are independent.

For a portfolio with weights $\mathbf{w}$ and covariance matrix $\Sigma$,
\[
  \sigma_p^2 = \mathbf{w}^{\!\top}\Sigma\,\mathbf{w} ,
\]
a quadratic form in the matrix language of \S\ref{sec:matmul}.
}

\analogybreak{
The Sharpe ratio penalises upside and downside volatility identically, which no
investor does, and it depends on a variance that exists here only because
$\hat\alpha > 2$. Had the tail index fallen below 2, the denominator would be
infinite and the ratio undefined. It is also easy to inflate: a strategy selling
deep out-of-the-money options shows a superb Sharpe for years and then loses
everything, because the risk it carries is entirely in the tail the denominator
cannot see.
}

\whyhere{
No Sharpe ratio is reported in this thesis, and the reason is worth stating: the
generators are evaluated on distributional and temporal fidelity, not on trading
performance. A synthetic series with a plausible Sharpe but wrong tail behaviour
would be worse than useless for risk work, which is the intended application. The
portfolio mathematics above is included to mark the boundary --- these are
univariate models, and $\Sigma$ is exactly what they do not provide.
}

%% ─────────────────────────────────────────────────────────────────────────────
\section{Structural Breaks and Regime Change}
\label{sec:breaks}

\situation{
A river's flood records span eighty years. Halfway through, a dam was built
upstream. Averaging all eighty years describes neither the pre-dam river nor the
post-dam one.
}

\problem{
Every method in this book has assumed stationarity --- that the rules generating
the data do not change (\S\ref{sec:stationarity}). Real markets break that
assumption, and this project's dataset contains a textbook case.
}

\workedarith{
\textbf{MOEX, February 2022.} The Russian invasion of Ukraine produced, in order:

\medskip
\begin{tabular}{@{}ll@{}}
\toprule
Date & Event \\
\midrule
2022-02-24 & single-day return of $-33.3\%$ \\
2022-02-25 to 2022-03-24 & exchange suspended, 27 trading days absent \\
after & trading resumes under capital controls \\
\bottomrule
\end{tabular}

\medskip
\textbf{What the single day does to the statistics.} MOEX's sample kurtosis is
$65.70$, the largest of the five markets, and it is driven overwhelmingly by this
one observation. Per \S\ref{sec:kurtosis} that number estimates nothing --- the
population kurtosis is infinite --- but it does show how one day dominates a
20-year sample.

\medskip
\textbf{What the suspension does.} Twenty-seven trading days are simply missing.
Any calculation treating consecutive rows as consecutive days silently bridges a
month-long gap. The return computed across the reopening is not a daily return at
all; it is a month's worth of accumulated news compressed into one row.

\medskip
\textbf{How the pipeline handles it.} The closure is recorded explicitly as a
known gap rather than being interpolated or discarded silently. The alternatives
are worse: interpolating invents data during precisely the period when the true
value was undefined, and dropping the market entirely would remove the most
informative regime change in the dataset.

\medskip
\textbf{What it means for the walk-forward protocol.} Folds are contiguous blocks
in time, so one fold contains the break and its neighbours do not. That fold's
training and test distributions genuinely differ, and a model that scores badly
there may be behaving correctly --- it was trained on one regime and evaluated on
another.
}

\formaldef{
A \textbf{structural break} is a change in the data-generating process at some
point in the sample: a shift in mean, variance, or dependence structure. It
violates weak stationarity (\S\ref{sec:stationarity}), since the moments are no
longer constant in time.

Tests exist for a break at a known date (Chow) and at an unknown one
(Quandt--Andrews, Bai--Perron). A \textbf{regime-switching} model instead treats
the state as latent and estimates transition probabilities.
}

\analogybreak{
Structural breaks are not outliers to be cleaned. Removing 2022-02-24 from MOEX
would produce a better-behaved series describing a market that does not exist ---
one in which invasions do not happen. The extreme observations are the phenomenon
the thesis exists to model, and a generator that cannot produce them has failed at
its actual job.

It does not follow that anything goes. There is a real difference between a
genuine $-33.3\%$ move and a data error, and only domain knowledge distinguishes
them.
}

\whyhere{
The MOEX case is why the train/validation/test split is strictly temporal
(\S\ref{sec:stochproc}) and why results are reported per market rather than
pooled for the headline comparison. It is also the honest limit on every
stationarity-dependent metric in Chapter~\ref{ch:evaluation}: they assume the test
distribution matches the training one, and for MOEX across 2022 that is false by
construction. Where a model does poorly on MOEX, the regime change is a candidate
explanation that must be ruled out before architecture is blamed.
}
